% This file is the shared source for the browser and PDF blueprints.

\chapter{Purpose and current position}

Zenith is an effect library written in Lean. Its formalization does not try to
prove every runtime feature at once. It formalizes one small boundary, proves
the model, relates the model to a production representation, and then extends
the boundary. This blueprint is the single status map for that work.

\section{How to read this blueprint}

Each completed item has a Lean declaration name. A green check means that the
named declaration exists and Lean's kernel checked it. The browser buttons
become public source links after an optional API-documentation deployment.
This does \emph{not} automatically mean that all production behavior is
proved. The evidence label states the exact scope:

\begin{description}
  \item[Kernel-proved] A theorem in \texttt{Prop} holds for every value in its
    stated Lean domain.
  \item[Production-connected] A checked relation connects an abstract model
    to an existing Zenith representation.
  \item[Fixture-checked] A regression program must elaborate, or must fail to
    elaborate. This checks representative elaborator behavior.
  \item[Planned] The property is a design target. It is not yet a theorem or a
    complete production connection.
\end{description}

The browser version also draws the dependencies declared with \texttt{uses}.
The formalization study guide gives the detailed method and reading order.

\section{Status board}

\begin{center}
\begin{tabular}{p{0.27\textwidth}p{0.18\textwidth}p{0.45\textwidth}}
\textbf{Track} & \textbf{Status} & \textbf{Current boundary} \\
\hline
Type algebra & Complete & Requirements as intersections and errors as unions,
  including canonical normal forms. \\
Production type connection & Complete with stated limits & Stable service rows,
  nested error shapes, and variance connect to production values. Arbitrary
  Lean error-type canonicalization remains fixture-checked. \\
Sequential semantics & Complete & A pure evaluator, a typed stack machine,
  and finite-step refinement cover six sequential instructions. \\
Executable sequential routing & Complete & The extracted dispatcher has
  reduction and conformance theorems for instruction and continuation routing.
  \\
Interpreter driver and async protocol & Active next boundary & Preserve the
  same-task and immediate-callback fast paths while specifying one resume-gate
  protocol or a one-loop driver. \\
Whole-runtime correctness & Planned & Fibers, interruption, scheduling,
  external IO, logging, diagrams, clocks, promises, references, and queues
  are outside the present theorem. \\
\end{tabular}
\end{center}

\section{Method: model, connect, refine}

The formalization has four repeated stages:

\begin{enumerate}
  \item Select a behavior small enough to state without unrelated runtime
    details.
  \item Define a semantic model and prove its internal laws.
  \item State a correspondence between the model and a production
    representation.
  \item Prove that a production transition simulates the model transition,
    then link that transition to executable code when possible.
\end{enumerate}

This sequence avoids a common error: treating a theorem about a convenient
model as a theorem about the executable interpreter. Each new connection
reduces that gap, but the blueprint keeps the remaining gap visible.

\chapter{Type-level effect algebra}

\section{Abstract requirements and errors}

\begin{definition}[Requirement intersection]
\label{def:requirement-intersection}
\lean{Zenith.Formalization.TypeAlgebra.Requirement.and_isGreatestLowerBound}
\leanok
For a service type, \texttt{Requirement} has \texttt{any}, \texttt{service},
and \texttt{and} forms. Its subtyping order uses reverse service-set
inclusion. Therefore \texttt{and} is a greatest lower bound: a requirement
for both services is below either requirement.
\end{definition}

\begin{theorem}[Requirement canonical forms]
\label{thm:requirement-normal-forms}
\uses{def:requirement-intersection}
\lean{Zenith.Formalization.TypeAlgebra.Requirement.normalForm_eq_of_equivalent}
\leanok
For ordered service leaves, semantically equivalent requirements have exactly
equal sorted, duplicate-free normal forms. This moves associativity,
commutativity, and idempotence from equivalence to exact equality of the
canonical representation.
\end{theorem}

\begin{definition}[Error union]
\label{def:error-union}
\lean{Zenith.Formalization.TypeAlgebra.ErrorType.or_isLeastUpperBound}
\leanok
For a failure type, \texttt{ErrorType} has \texttt{nothing}, \texttt{failure},
and \texttt{or} forms. Its subtyping order uses ordinary failure-set
inclusion. Therefore \texttt{or} is a least upper bound.
\end{definition}

\begin{theorem}[Error canonical forms]
\label{thm:error-normal-forms}
\uses{def:error-union}
\lean{Zenith.Formalization.TypeAlgebra.ErrorType.normalForm_eq_of_equivalent}
\leanok
For ordered failure leaves, semantically equivalent errors have exactly equal
canonical normal forms.
\end{theorem}

\section{Connections to Zenith values}

\begin{theorem}[Service-row provision]
\label{thm:service-row-provision}
\uses{def:requirement-intersection,thm:requirement-normal-forms}
\lean{Zenith.Formalization.ServiceRows.canProvide_provides, Zenith.Formalization.ServiceRows.canProvide_nonempty_of_provides}
\leanok
Typed \texttt{Environment.CanProvide} evidence implies the abstract
provision relation for stable keyed service rows. Conversely, provision plus
row coherence gives nonempty typed projection evidence. The coherence
condition prevents incompatible entries with the same key from being treated
as one service.
\end{theorem}

\begin{theorem}[Nested error-shape injection]
\label{thm:error-shape-injection}
\uses{def:error-union,thm:error-normal-forms}
\lean{Zenith.Formalization.ErrorShape.joinUpperBound}
\leanok
The selected production representation is nested \texttt{Sum}. A checked
syntax describes its shapes and proves branch-injection evidence for their
join. The automatic flattening, ordering, and duplicate removal of arbitrary
Lean error types remains fixture-checked because arbitrary \texttt{Type}
values have no stable public error key.
\end{theorem}

\begin{theorem}[Effect variance and composition]
\label{thm:effect-variance}
\uses{thm:service-row-provision,thm:error-shape-injection}
\lean{Zenith.Formalization.VarianceLaws.flatMapMeetJoin}
\leanok
The production API checks the ZIO-shaped rule: environment requirements meet,
while error channels join during heterogeneous sequential composition. The
focused coercions and adaptation functions are checked by the accompanying
variance laws and compile-time tests.
\end{theorem}

\chapter{Sequential interpreter correctness}

\section{Semantic model}

\begin{definition}[Deterministic pure evaluation]
\label{def:sequential-evaluation}
\lean{Zenith.Formalization.SequentialCore.evaluates_deterministic}
\leanok
\texttt{SequentialCore.Program} contains the first six sequential
instructions: completed exits, \texttt{flatMap}, \texttt{foldCauseM},
environment adaptation, environment reads, and supplied environments.
\texttt{Evaluates} gives their big-step meaning and has deterministic final
exits.
\end{definition}

\begin{theorem}[Evaluation reaches the typed machine]
\label{thm:evaluation-machine}
\uses{def:sequential-evaluation}
\lean{Zenith.Formalization.SequentialMachine.evaluation_runs_to_halt}
\leanok
Every direct evaluation of the model reaches the same final exit through a
typed continuation-stack machine. This establishes the operational model for
the next refinement steps.
\end{theorem}

\begin{theorem}[Production stack correspondence]
\label{thm:stack-correspondence}
\uses{thm:evaluation-machine}
\lean{Zenith.Formalization.SequentialRuntimeStack.corresponding_size}
\leanok
Every pure continuation stack has a corresponding production \texttt{Stack},
and the correspondence preserves the frame count exactly. It also records the
saved environments and projection evidence that the runtime stack carries.
\end{theorem}

\section{Production-shaped transitions}

\begin{theorem}[One production transition refines one machine transition]
\label{thm:sequential-step-refinement}
\uses{thm:evaluation-machine,thm:stack-correspondence}
\lean{Zenith.Formalization.SequentialRuntime.step_refines}
\leanok
For each model-machine step, a corresponding pure production transition
exists. The relation uses lowered production \texttt{ZCore} instructions,
the real production continuation stack, and saved environment evidence.
\end{theorem}

\begin{theorem}[Finite sequential refinement]
\label{thm:sequential-steps-refinement}
\uses{thm:sequential-step-refinement}
\lean{Zenith.Formalization.SequentialRuntime.steps_refine}
\leanok
The one-step simulation extends to every finite execution of the selected
sequential model.
\end{theorem}

\section{Executable dispatcher connection}

\begin{theorem}[Executable instruction routing conforms]
\label{thm:dispatcher-run}
\uses{thm:sequential-step-refinement}
\lean{Zenith.Formalization.SequentialDispatcher.run_models_step}
\leanok
\texttt{Z.Runtime.Sequential.run} is the extracted executable dispatcher.
For each lowered model instruction, its routing result corresponds to one
production transition.
\end{theorem}

\begin{theorem}[Executable continuation routing conforms]
\label{thm:dispatcher-resume}
\uses{thm:dispatcher-run,thm:stack-correspondence}
\lean{Zenith.Formalization.SequentialDispatcher.success_models_step, Zenith.Formalization.SequentialDispatcher.failure_models_step}
\leanok
The executable success and failure routing functions make the same choices as
the production transition relation for lowered continuation frames.
\end{theorem}

\chapter{Next proof boundaries}

\section{Interpreter driver and asynchronous registration}

\begin{definition}[Resume-gate protocol]
\label{def:resume-gate}
\uses{thm:dispatcher-run,thm:dispatcher-resume}
\notready
The next recommended boundary specifies one asynchronous registration and
resume-gate protocol. It must distinguish completion during registration from
completion after registration. The former must continue on the current task;
the latter may need a task hop. This is a planned model and proof, not a
completed theorem.
\end{definition}

\begin{definition}[One-loop interpreter driver]
\label{def:interpreter-driver}
\uses{thm:dispatcher-run,thm:dispatcher-resume}
\notready
An alternative next boundary extracts a small driver around the present
dispatcher. The goal is to make the decision part total where possible. The
complete driver can still need \texttt{partial}, because an effect can wait
for a callback or run forever. This item needs a design decision before work
starts.
\end{definition}

\section{Runtime work outside the current theorem}

\begin{definition}[Whole-runtime extension]
\label{def:whole-runtime-extension}
\uses{def:resume-gate,def:interpreter-driver}
\notready
Later boundaries can cover interruption, fibers, scheduling, external IO,
logging, diagrams, clocks, promises, mutable references, queues, streams, and
scopes. They should enter one behavior at a time. They are not implied by the
current sequential proofs.
\end{definition}

\section{Performance obligations}

Formal proofs in \texttt{Prop} are erased, but an interpreter refactor can
still add allocations, closures, queue operations, or task creation. Every
driver or async change must preserve these measured behavior requirements:

\begin{enumerate}
  \item Ordinary synchronous steps continue on the caller task.
  \item A callback that completes during registration resumes inline.
  \item Sequential steps do not pass through a queue.
  \item The synchronous and immediate-async benchmark paths do not gain task
    creation.
  \item The benchmark comparison covers \texttt{flatMap}, \texttt{sync},
    error recovery, immediate async completion, and fork/join.
\end{enumerate}

These are engineering acceptance conditions, not kernel theorems. The exact
baseline and commands are in \texttt{docs/interpreter-refactor-plan.md} and
\texttt{Benchmarks/README.md}.

\chapter{Maintenance rule}

Update this blueprint in the same change as a proof-boundary change. A new
theorem needs a declaration link, dependencies, a status label, and a clear
statement of what remains outside its domain. Run the checked-formalization
build and the declaration-check file before claiming that a status item is
complete.
